\editionPage[\pdfbookmark[0]{Fiches outils}{bm-outils-debut-parent}]{outils-debut}{Fiche outil}{Poser les quatre opérations}{À consulter au fil des chapitres.}
\begin{blocEdition}{outils-debut}
\editionEtape{Aligner}Unités sous unités, dizaines sous dizaines ; pour les décimaux, virgules alignées et rangs vides complétés par des zéros.\editionEtape{Retenues}Dix unités s’échangent contre une dizaine, dix centièmes contre un dixième. Écrire les retenues petites et au bon rang.\editionEtape{Division}Encadrer le dividende par des multiples du diviseur, abaisser chiffre à chiffre, écrire 0 au quotient quand le chiffre abaissé ne suffit pas. Le reste est toujours plus petit que le diviseur.\editionEtape{Contrôler}Avant : estimer avec un ordre de grandeur. Après : vérifier la soustraction par une addition, la division par quotient fois diviseur plus reste.\editionRetenir{Une opération posée se termine par son contrôle. Un résultat loin de l’estimation se refait.}
\end{blocEdition}
\editionPage{outils-2}{Fiche outil}{Résoudre un problème en quatre étapes}{À consulter au fil des chapitres.}
\begin{blocEdition}{outils-2}
\editionEtape{Comprendre}Raconter l’histoire du problème avec ses propres mots : ce que l’on sait, ce que l’on cherche. Un mot comme « plus » ne décide jamais de l’opération.\editionEtape{Modéliser}Choisir les opérations, au besoin avec un schéma en barres : le tout en haut, les parties en dessous ; pour comparer, deux barres alignées et l’écart.\editionEtape{Calculer}Calcul mental quand c’est raisonnable, calcul posé sinon. Garder les étapes visibles, une ligne par calcul.\editionEtape{Répondre}Écrire une phrase avec l’unité, puis contrôler la vraisemblance : un reste plus grand que le tout ou une voiture de 50 m signalent une erreur.\editionRetenir{Comprendre, modéliser, calculer, répondre : quand on bloque, chercher laquelle des quatre étapes coince.}
\end{blocEdition}
\editionPage{outils-3}{Fiche outil}{Construire proprement en géométrie}{À consulter au fil des chapitres.}
\begin{blocEdition}{outils-3}
\editionEtape{Préparer}Lire tout le programme de construction, faire un croquis à main levée codé, repérer par quoi commencer : souvent le segment le plus long.\editionEtape{Choisir l’instrument}Règle pour tracer et mesurer, équerre pour les angles droits, compas pour reporter une longueur ou tracer un cercle. Un crayon fin et taillé change tout.\editionEtape{Coder}Petits traits pour les longueurs égales, carré pour l’angle droit : les codes disent les propriétés voulues, même quand le tracé est imparfait.\editionEtape{Vérifier}Mesurer les longueurs demandées, tester les angles droits à l’équerre, contrôler que les points nommés sont tous placés.\editionRetenir{Garder les traits de construction : arcs de compas et lignes légères racontent la méthode.}
\end{blocEdition}
\editionPage{outils-4}{Fiche outil}{Lire un tableau ou un graphique}{À consulter au fil des chapitres.}
\begin{blocEdition}{outils-4}
\editionEtape{Observer d’abord}Titre, grandeurs, unités : que raconte ce document ? Dans un tableau à double entrée, une case se lit au croisement d’une ligne et d’une colonne.\editionEtape{Lire les axes}Repérer le pas des graduations : un carreau peut valoir 1, 2, 5 ou 100. Sur une courbe, l’abscisse se lit avant l’ordonnée.\editionEtape{Croiser}Vérifier les totaux d’un tableau d’enquête ; comparer des barres par leurs hauteurs lues, pas à l’œil ; chercher les points les plus hauts et les plus bas d’une courbe.\editionEtape{Rester prudent}Entre deux points de mesure, la courbe ne garantit rien ; un graphique sans unités ne prouve rien.\editionRetenir{Un graphique répond vite à « combien » et « quand » ; le tableau garde les valeurs exactes. Les deux se contrôlent l’un l’autre.}
\end{blocEdition}
